
%The trivial threshold signature scheme
%represents the public key as the set of $n$ public keys for each party.
%To sign,
%each participant produces a signature individually;
%the output signature is simply the concatination of all $t$ signatures.
%Verification requires each signature to be individually verified with respect
%to its respective public key.

%However, the trivial threshold signature scheme is undesirable for two reasons.
%First, the performance overhead of the trivial threshold scheme scales linearly
%relative to $t$.
%Second,
%the verification algorithm of the trivial threshold scheme is not compatible with existing single-party threshold signature schemes.

In this chapter,
we focus our attention on deterministic threshold Schnorr signatures;
i.e,
a deterministic threshold signature scheme that is compatible with the (single-party) Schnorr signature verification algorithm~\cite{Schnorr91}.
As a reminder,
threshold signature schemes allow a subset of at least $\minsigners$ out of $n$ possible parties to cooperate to produce a signature over a single message,
while preserving security when up to $t-1$ signers may be corrupted,
where $\minsigners \leq n$ and  $t \leq \minsigners$.


Efficient (two-round) threshold Schnorr signatures exist in the randomized setting~\cite{KomloG20,BellareCKMTZ22,RuffingRJSS22},
where each party is assumed to securely maintain state between signing rounds and have access to good sources of randomness.
However,
efficient and \emph{deterministic} threshold Schnorr signatures has thus far remained an elusive goal.

\paragraph{\bf Why Deterministic Threshold Signatures?}
Producing threshold signatures in a deterministic manner is useful for two reasons.
%
First,
it is useful as a general defense-in-depth measure,
to protect against the event of temporarily losing access to good sources of   randomness~\cite{RistenpartY10},
such as if a machine randomly rebooted.

Second,
threshold Schnorr signature schemes generally require participants to perform multiple rounds of communication before a joint signature can be issued.
As such, participants must keep state between each round,
and carefully delete state once a signing protocol is completed.
In the setting where the signer must produce signatures at increasing scale and
in a concurrent setting,
managing state can become a significant performance bottleneck.
Further,
%because in practice machines can randomly go offline at any time,
secure management of secret state can be considered a security risk.
In particular,
for threshold Schnorr signatures,
participants must generate secret nonces for each signing session.
If a participant's state is mismanaged in such a way that it is used even twice,
a fatal key-recovery attack is possible.
However, in the setting where machines can go offline at any time,
and signing is done at scale and concurrently,
such careful management of secret state becomes even more challenging.
%\ian{I'm not doing a careful read over the intro, as requested, but I wanted to leave a note for later that the checkpointing-and-restoring-VMs problem is also a good example of where deterministic signatures are useful.}


\paragraph{\bf The Challenge of Efficient Deterministic Multi-Party Schnorr Signatures.}
While the goal of efficient and deterministic multi-party Schnorr signature schemes is desirable,
producing such schemes has proven difficult.
The challenge can be summarized as follows.
While honest participants can certainly  pick their nonces deterministically
(say, by hashing their secret signing share and the message),
a malicious party might deviate from the protocol and pick its nonce at random.
Recall that for Schnorr signatures,
the output signature $\sigma = (R,z)$ is a commitment $R$ and response $z$,
satisfying the relation $g^z = R \cdot \pk^c$ for a challenge $c \gets H(R,\pk,m)$.
In the setting for threshold Schnorr signatures,
where $(R,z)$ are contributed to by all signing parties,
one party deviating from the protocol results in the challenge $c$ likewise changing.
If honest parties cannot verify that other parties followed the protocol,
such a deviation would allow an adversary to perform a key recovery attack with as little as two signing queries.

Prior deterministic multi-party Schnorr signature schemes in the literature~\cite{NickRSW20,GarillotKMN21,KondiOR23}
can be broken down into two general approaches.
We give an overview of the approaches that support the general threshold setting with general $(t, \minsigners, n)$ in Table~\ref{table:efficiency},
with more context now.

The first approach requires that signers output zero-knowledge proofs that a Pseudorandom Function (PRF) was evaluated correctly~\cite{NickRSW20,GarillotKMN21}.
However,
this approach adds undesirable performance and complexity overhead.
For example,
while MuSig-DN~\cite{NickRSW20} requires only two network rounds,
proof generation requires approximately $1$~second (regardless of the number of signers) for $256$-bit security,
and verification requires at least $10n$\,ms.
%\ian{That's not what their Table 1 says? Where did this number come from?}
%\chelsea{edited; this roughly follows that they report in their table}
%\ian{For Musig-DN, $\minsigners = t$, right? So perhaps using $t$ is better, particularly because ``$100\mu$\,ms'' looks like you're mixing microseconds and milliseconds?}
%\chelsea{it is technically n, i think saying t is a little too generoug}
Similarly,
the scheme by Garillot et al.~\cite{GarillotKMN21}
requires three rounds of communication between signers and high complexity and bandwidth overhead.

The second approach  by Kondi et al.~\cite{KondiOR23} is to use as a black box a Pseudorandom Correlation Generator (PCG)~\cite{BoyleCGIKS20}  to generate nonces in a verifiable manner~\cite{KondiOR23}.
However, while the PCG used in their construction can support $2$-of-$2$ signing,~\cite{OrlandiSY21,KondiOR23},
it remains an open problem as to how to extend such a PCG to the general threshold setting for any $t$, $\minsigners$,  or $n$~\cite{KPC24}.
As such,
the scheme by Kondi et al.~\cite{KondiOR23} supports only the $t=\minsigners=n=2$ setting.
%This is because the PCG requires the verifier to maintain a secret verification key unknown to the prover;
%otherwise all security is lost.

Therefore
in this chapter,
we address the following question:

\begin{quote}
  \it
Can we design a simple, deterministic, and stateless two-round threshold Schnorr signature scheme,
with similar efficiency to existing randomized schemes,
and acceptable security assumptions in practice?
\end{quote}

\subsection{Our Results}
We answer the above question in the affirmative.
In particular, we present \glitter,
a deterministic and stateless threshold Schnorr signature scheme.
For moderately sized groups (i.e., when $n \leq 20$),
\glitter is more than an order of magnitude more efficient than existing deterministic schemes in the literature;
it is three orders of magnitude more efficient when $n \leq 10$.
For even larger groups,
\glitter scales linearly relative to the number of processing cores available to most machines.
Signers in \glitter are required to generate and maintain secret keys,
but otherwise do not need to manage secret state.

To achieve efficient determinism,
\glitter requires three tradeoffs.
The first tradeoff is in the number of required signers;
\glitter requires $\minsigners \geq 2t-1$ parties to participate in signing.
This requirement is because \glitter assumes the honest majority model.
We discuss in Section~\ref{hon-maj-obs} why the assumption of an honest majority is acceptable for some real-world applications.

The second tradeoff is in the required overhead as the signing set grows large.
Under the hood, \glitter employs a replicated secret sharing scheme~\cite{MitsuruSN89},
and so requires participants to store a secret key of size
${n-1  \choose t-1}$.
As such,
\glitter outperforms other schemes in the literature for moderately sized groups,
but incurs a crossover point as the signing set grows large.

The final tradeoff is that \glitter as defined does not support identifiable abort.
While the protocol allows honest participants to identify misbehavior has occurred,
the protocol does not support identifying which participant caused in the abort
(which is possible for MuSig-DN~\cite{NickRSW20}, for example).
However,
we give a discussion on how \glitter can be extended to support identifiable abort,
and indeed robustness,
in the \emph{honest supermajority} setting (i.e., when $\minsigners \geq 3t-2$).

\begin{table}[t]
	\centering
%		\resizebox{\textwidth}{!}{
		\begin{tabular}{ c | c | c | c | c | c  }
		 \hline
      %& \rot{Deterministic?} & \rot{Stateless?}
      & \rot{Computation}
      & \rot{Bandwidth}
      & \rot{Num. Rounds} & \rot{Assumptions} &
      \rot{Min. Signers }   \\
             \hline
             % TODO
      %FROST \cite{KomloG20,BellareCKMTZ22} & \xmark & \xmark & 0.1  & $t$ group & 2 & AOMDL & t  & n \\
             % TODO add back in for the final version
      %\hline
      MuSig-DN~\cite{NickRSW20} & $14210 + 24 (n - 1)$ group & 1189~bytes & 2 & DDH & $n$  \\
      \hline
      \multirow{2}{*}{ GKMN21~\cite{GarillotKMN21}} & $132,000(t-1)$ AES & \multirow{2}{*}{ $1.01(t-1)$~MB}
      & \multirow{2}{*}{3} & \multirow{2}{*}{PRF} &  \multirow{2}{*}{$t$}  \\
      &   $256(t-1)$ field &  & &   \\
      %\hline
      %\multirow{2}{*}{KOR23~\cite{KondiOR23}} & \multirow{2}{*}{\xmark} & \multirow{2}{*}{\cmark} & \multirow{2}{*}{\cmark} &
      %\multirow{2}{*}{0.5} & todo & \multirow{2}{*}{2} & \multirow{2}{*}{DCR+Strong RSA} & \multirow{2}{*}{2} & \multirow{2}{*}{2} \\
      \hline
      \glitter  & $2t^2 -t + 2$ group &
      \multirow{2}{*}{97~bytes} & \multirow{2}{*}{2} & \multirow{2}{*}{DL} & \multirow{2}{*}{$2t-1$}   \\
      (this work) &  $2{n-1 \choose t-1}$ field,  $2{n-1 \choose t-1}$ hash &  & &   \\
      \hline
	\end{tabular}
%}
  \medskip
	\caption{ Comparison of multi-party deterministic Schnorr signature schemes
  in the ROM that support general choices of $n$.
  }
\label{table:efficiency}
\end{table}


\paragraph{\bf Verifiable Pseudorandom Secret Sharing (VPSS).}
As a building block,
we first formalize a cryptographic primitive that we call a \emph{Verifiable Pseudorandom Secret Sharing (VPSS)} scheme.
While pseudorandom secret sharing is a standard notion in the literature,
and verifiability of such schemes had been employed implicitly in maliciously secure MPC-based schemes~\cite{BenhamoudaBGHIN21},
we present a formal definition for a VPSS and
%\ian{Were you going to add something like ``following Bellare et al.~\cite{something},''?}
%\chelsea{I don't think it is necessary to elaborate on it, beyond a citaiton}
give game-based security notions~\cite{BellareR06},
building on existing notions of verifiable random functions in the literature~\cite{MicaliRV99,NaorPR99,Dodis03,GalindoLOW21}.
%taking inspiration from prior security notions for distributed VRFs~\cite{MicaliRV99,NaorPR99,Dodis03,GalindoLOW21}.

Intuitively, a pseudorandom secret sharing scheme allows a coalition of players holding pre-established secret shares of a random secret to generate shares of additional pseudorandom secrets.
%We additionally require that a pseuodorandom secret sharing scheme be \emph{publicly verifiable};
%i.e., it is possible to determine that the output held by each participant is valid,
%simply by verifying commitments to secret shares output by each participant.
%
A \emph{verifiable pseudoradnom secret sharing scheme} is simply an extension
to allow for public verifiability by collectively verifying the outputs for the set of participants.
As such, misbehavior of players in a VPSS can be detected assuming some threshold of honest participants.
As we will see in the case of \glitter,
a VPSS allows for more efficiently verifying the correctness of the output of a pseudorandom function that is distributed among a set of players,
without requiring each player individually to produce a zero-knowledge proof that it followed the protocol.

We then define a concrete VPSS construction that we call \DVRFOne
that builds upon the pseudorandom secret sharing scheme by Cramer et al.~\cite{CramerDI05}.
We augment the scheme by Cramer et al.\ to additionally define a verification algorithm that is publicly verifiable even when players only publish commitments to their shares;
i.e., it preserves secrecy of players' shares while allowing players to verify that all other players followed the protocol honestly.
\DVRFOne is efficient for moderately sized groups;
concretely, evaluation requires ${n-1\choose t-1}$ sub-microsecond
hash and field operations,
and verification requires $2t^2 - t$ group exponentiations,
where $t$ is the corruption threshold.
We prove that \DVRFOne is secure assuming the existence of a secure hash function in the honest majority model,
i,e.,
when $\minparticipants \geq 2t-1$,
where $\minparticipants$ is the minimum number of participants required to participate in the protocol.
Additionally, we require that as $n$ grows, $t$ remains $\BigO(1)$ so that
${n-1\choose t-1} = \text{poly}(n)$.


\paragraph{\bf A New Two-Round, Deterministic, and Stateless Threshold Schnorr Signature.}
Next,
we introduce \glitter,
a new two-round, deterministic, and stateless threshold Schnorr signature.
\glitter is an order of magnitude more efficient than related schemes in the literature for moderately sized groups ($n \le 20$),
due to its use of \DVRFOne as a building block.
Furthermore,
\glitter supports the general threshold setting for any $t, \minsigners$ and $n$,
so long as $\minsigners \geq 2t-1$ and $\minsigners \leq n$.

At a high level,
\glitter uses \DVRFOne in the first round of signing to generate participant nonces and commitments;
participants publish their commitments without having to maintain state of their (secret) nonce.
Then, in the second round,
participants use \DVRFOne to re-derive their nonce and commitment,
and verify all other participants' commitments,
ensuring that other participants followed the protocol correctly.
If the verification check holds,
participants derive the group commitment as the aggregation of all participants' individual commitments,
and generate a signature share as a combination of their nonce, challenge, and secret signing share.
The output signature is the aggregated group commitment and an aggregation of all participants' signature shares,
and can be verified using the single-party Schnorr verification algorithm.,

We prove the unforgeability of \glitter assuming the hardness of the discrete logarithm problem in the random oracle model,
assuming $\minsigners \geq 2t-1$,
the adversary corrupts at most $t-1$ players,
and ${n-1\choose t-1} = \text{poly}(n)$.
We describe in Section~\ref{hon-maj-obs} why requiring honest majority assumptions can be practical for some real-world deployments of threshold signatures,
at the benefit of improved simplicity and performance.


\paragraph{\bf Comparison with Related Schemes.}
We give an overview of \glitter in Table~\ref{table:efficiency} in comparison to related
schemes in the literature.
In the table,
we review only schemes that support a general number $n$ of parties (i.e., we exclude $2$-of-$2$ schemes).
Computation is given with respect to the operations that dominate;
estimates for MuSig-DN are given with respect to a $256$-bit elliptic curve;
more information is given in Section~\ref{performance-estimates}.
Bandwidth similarly is given with respect to a $256$-bit elliptic curve~\cite{NickRSW20,GarillotKMN21}.

\paragraph{\bf Performance Analysis of \glitter.}
To estimate the practicality of \glitter for various choices of $n, t$, and $\minsigners$,
we implement the scheme and provide concrete benchmarks.

\glitter is highly efficient for moderately sized groups,
and so is a good choice for such settings.
As we show in greater detail in Section~\ref{sec:performance},
for a group where $2t-1 \leq \minparticipants \le n \leq 20$,
\glitter signing operations require less than 100 milliseconds of computation for each signer, and for $n \le 10$, less than 1 millisecond.
%\ian{Wait; that's not what the figure shows?  Where did this number come from?}
%\chelsea{can you update if this is incorrect? this looks correct to me}
%\ian{For $n \le 20$, the figure shows 100 milliseconds, not 100 microseconds?}
Signing however increases in cost relative to the corruption threshold;
for signing sets of size $n=25, t=11$, the
computational overhead for each signer on a single core requires one second.
However, \glitter is highly parallelizable, showing almost perfect linear speedups for up to 32 cores for that size of signing set.

For comparison,
MuSig-DN requires about $1$~second in computational overhead per signer~\cite{NickRSW20},
regardless of the size of the signing set.
As such, \glitter is a practical choice for moderately sized groups of $n \leq 25$,
whereas MuSig-DN or GKMN21~\cite{GarillotKMN21} may be good choices as the set of signers grows large,
depending on the parallelization of each scheme.

\paragraph{\bf Our Contributions.}
In summary,
the contributions of our work are as follows:
\begin{itemize}
  \item We formalize the definition of a Verifiable Pseudorandom Secret Sharing (VPSS) scheme,
    and give game-based security notions.
  \item We define a concrete VPSS scheme that we call \DVRFOne,
    which extends the pseudorandom scheme by Cramer et al.~\cite{CramerDI05} to allow for public verifiability,
    assuming an honest majority of participants.
  \item We then present \glitter,
    a deterministic and stateless two-round threshold Schnorr signature scheme.
    \glitter uses \DVRFOne as a building block to generate participant nonces and commitments,
    and verify that participants performed this action correctly.
  \item We prove that \glitter is secure under the discrete logarithm problem
    in the random oracle model,
    assuming \DVRFOne is a secure VPSS,
    the adversary corrupts fewer than $t$ parties,
    at least $\minsigners \geq 2t-1$ parties participate in the signing protocol,
    and ${n-1\choose t-1} = \text{poly}(n)$.
  \item We provide performance benchmarks for \glitter for different sizes of signing sets,
    and show that parallelization enables a linear speedup in signer computation.
\end{itemize}

\subsection{Observations of Honest Majority Assumptions in Practice}
\label{hon-maj-obs}

While honest minority assumptions may be desirable for some real-world applications,
we observe that honest majority assumptions may be an acceptable tradeoff for other applications in exchange for statelessness, improved performance, and protocol simplicity.

For applications that can easily support additional signers,
moving from an honest minority setting to honest majority can be relatively straightforward.
For example,
an application that currently uses a $(t=2, \minsigners=2, n=3)$ configuration can instead move to a $(t=2, \minsigners=3, n=4)$ configuration.

In settings that require liveness,
honest majority requirements are already assumed,
to ensure usability of the shared secret key even when corrupt players refuse to participate.
For example,
applications such as cryptocurrency wallets often implicitly require honest majority assumptions,
to ensure that corrupt players cannot prevent use of funds by simply refusing to participate in signing operations.

Finally, some applications may see the requirement for additional signers as an acceptable tradeoff to mitigate the security risk of protocol complexity,
as protocol complexity increases the risk of exploitable implementation errors~\cite{MakriyannisYG23}.

%Prior works on deterministic threshold Schnorr signatures have focused on
%schemes that are secure assuming an honest minority of signers~\cite{},
%at the tradeoff of performance.
%However, honest majority assumptions are implicitly required by implementations
%of threshold signature schemes in practice.
%In particular,
%while honest majority assumptions may not necessarily be required for
%unforgeability,
%these assumptions are strictly required for liveness.
%In particular,
%real-world implementations of threshold signatures require a sufficient number
%of honest players such that even if corrupt players refuse to participate in
%the protocol,
%signatures can still be produced,
%otherwise the secret key is rendered unrecoverable.

%As such,
%the tradeoff of honest majority assumptions for improved performance for moderately sized groups and simplicity may be a natural fit for some applications of threshold signature schemes in practice.








